\section{Hollywood Action Flicks}


In his inimitable way, \textbf{Gilbert Chesterton}, in a discussion about adventure or romance novels NOTE: not meant in the contemporary sense , anticipates the plots and defects of Hollywood action movies. The title of the essay is ``The Position of Sir Walter Scott''; in this segment he compares Scott to later adventure novel authors.

\begin{quotex}
It is in this quality of what may be called spiritual adventurousness that Scott stands at so different an elevation to the whole of the contemporary crop of romancers who have followed the leadership of Dumas.
\end{quotex}


In contrast this is how he describes Scott's inferiors:

\begin{quotex}
There has, indeed, been a great and inspiriting revival of romance in our time, but it is partly frustrated in almost every case by this rooted conception that romance consists in the vast multiplication of incidents and the violent acceleration of narrative. The heroes of Mr. Stanley Weyman scarcely ever have their swords out of their hands; the deeper presence of romance is far better felt when the sword is at the hip ready for innumerable adventures too terrible to be pictured. The Stanley Weyman 1855-1928, adventure novel writer hero has scarcely time to eat his supper except in the act of leaping from a window or whilst his other hand is employed in lunging with a rapier. ... In short, Mr. Stanley Weyman is filled with the conviction that the sole essence of romance is to move with insatiable rapidity from incident to incident.
\end{quotex}


You see the parallels with contemporary action films? Nothing but chase scenes and fight scenes for two hours, with barely a plot intertwined: the later Die Hard movies, the incessant and repetitive battles of The Matrix, the faux fights of Kill Bill. And these are the better ones, of a list that can go on and on.

In contrast, this is what he says about Sir Walter Scott:

\begin{quotex}
In the truer romance of Scott there is more of the sentiment of ``Ho! Still delay, thou art so fair''; more of a certain patriarchal enjoyment of things as they are -- of the sword by the side and the wine cup in the hand. Romance, indeed, does not consist by any means so much in experiencing adventures as in being ready for them.
\end{quotex}


A criticism of Scott is that his descriptions of details is too monotonous; in one case about the details of armour and costume, GKC responds:

\begin{quotex}
The only thing to be said about the critic is that he had never been a little boy. ... Not being himself romantic, he could not understand that Scott valued the plume because it was a plume, and the dagger because it was a dagger. Like a child, he loved weapons with a manual materialistic love, as one loves the softness of fur or the coolness of marble. One of the profound philosophical truths which are almost confined to infants is this love of things, not for their use or origin, but for their own inherent characteristics, the child's love of the toughness or wood, the wetness of water, the magnificent soapiness of soap. So it was with Scott, who had so much of the child in him.
\end{quotex}


Yes, the half-educated man can only deal with his abstractions. But the boy prefers direct sensual experience. And, yes, we boys love pageantry, dressing up in costume, and especially our weapons.


\flrightit{Posted on 2010-04-21 by Cologero}

\begin{center}* * *\end{center}

\begin{footnotesize}
\begin{sffamily}
\texttt{rkirk on 2010-04-26 at 11:11 said:}

``The Stanley Weyman hero has scarcely time to eat his supper except in the act of leaping form a window or whilst his other hand is employed in lunging with a rapier.''\\ Chesterton anticipated not only the action genre, but also the question of when Jack Bauer sleeps, eats, or goes to the bathroom.

\hfill

\texttt{Cologero on 2010-04-26 at 22:15 said:}

Or perhaps jack Bauer is a follower of Plotinus and has contempt for such lowly human bodily functions.

\end{sffamily}
\end{footnotesize}

